\section{Učení s~učitelem}\label{sec:ai-supervised}

Při učení s~učitelem máme pro každý trénovací objekt k~dispozici jak vstupní znaky, tak správný výstup. Nechť
\[
    \mathbf{X}_1,\ldots,\mathbf{X}_p
\]
jsou vstupní statistické soubory a
\[
    \mathbf{Y}=(y_1,\ldots,y_n)
\]
je výstupní statistický soubor. \(i\)-tý objekt popisuje vektor
\[
    \mathbf{x}_i=(x_{i1},\ldots,x_{ip}),
\]
a~trénovací množinu tedy můžeme zapsat
\[
    D=
    \set{
        (\mathbf{x}_1,y_1),\ldots,
        (\mathbf{x}_n,y_n)
    }.
\]

Model \(h\) vytváří pro nový objekt \(\mathbf{x}\) předpověď \(h(\mathbf{x})\). Při trénování porovnáváme předpovědi na trénovacích objektech se známými hodnotami \(y_i\) a~upravujeme model tak, aby se zvolená chyba zmenšovala.

\begin{figure}[ht]
    \centering
    \begin{tikzpicture}[every node/.style={font=\normalsize}]
    \node[task,fill=blue!10,minimum width=34mm] (data) at (0,0)
        {\(D=\set{(\mathbf{x}_i,y_i)}\)};
    \node[task,fill=orange!15,minimum width=30mm] (learn) at (4,0)
        {učení modelu};
    \node[task,fill=green!12,minimum width=25mm] (model) at (8,0)
        {\(h\)};
    \node[task,fill=blue!10,minimum width=25mm] (new) at (4,-1.8)
        {\(\mathbf{x}\)};
    \node[task,fill=green!12,minimum width=31mm] (pred) at (8,-1.8)
        {\(h(\mathbf{x})\)};
    \draw[diredge] (data) -- (learn);
    \draw[diredge] (learn) -- (model);
    \draw[diredge] (new) -- (pred);
    \draw[diredge] (model) -- (pred);
\end{tikzpicture}

    \caption{Základní postup učení s~učitelem.}
    \label{fig:ai-supervised}
\end{figure}

Podle povahy výstupu rozlišujeme zejména dvě úlohy:
\begin{itemize}
    \item \textbf{regrese} předpovídá číselnou hodnotu, například cenu bytu, spotřebu energie nebo dobu jízdy;
    \item \textbf{klasifikace} vybírá z~konečné množiny tříd, například spam a~běžnou zprávu nebo jednu z~několika kategorií obrazu.
\end{itemize}

Dobře fungující model musí zachytit obecný vztah mezi vstupy a~výstupem. Příliš jednoduchý model důležitou strukturu nevystihne; příliš složitý model může místo ní zachytit šum. Volba třídy modelů proto určuje kompromis mezi schopností přizpůsobit se datům a~schopností zobecnit.

\notebox{Slovo „učitel“ neznamená, že model vede člověk krok za krokem. Učitelem je zde známý výstup \(y_i\), s~nímž lze porovnat předpověď modelu.}
